\documentclass[letterpaper,11pt,leqno]{article}
\usepackage{paper}
\bibliographystyle{paper}
\hypersetup{pdftitle={Belegarbeit Sebastian Schöneich - Vergleich von Python Webframeworks}}
\newcommand{\bib}{paper.bib}
\newcommand{\pdf}{figures.pdf}
\renewcommand{\contentsname}{Inhaltsangabe}
\renewcommand{\listfigurename}{Abbildungsverzeichnis}
\renewcommand{\listtablename}{Tabellenverzeichnis}

\begin{document}
\title{Vergleich von Python-Webframeworks\\Django vs. FastAPI vs. Flask}
\author{Sebastian Schöneich
\thanks{Duale Hochschule Sachsen - Standort Dresden, SARAD GmbH. Ich bedanke mich bei Komillitonen und Kollegen für Einblicke und
Kommentare zur Arbeit und zur Herangehensweise}}
\date{Januar 2026}   
\available{https://github.com/NeroAiur/Belegarbeit_03-PUBLIC}
\begin{titlepage}
\maketitle
Die Wahl des optimalen Webframeworks für ein Projekt ist entscheidend für viele Aspekte der
Softwarelösung zum Beispiel Effizienz, Skalierbarkeit und Funktionsumfang. Diese Belegarbeit
analysiert und vergleicht die drei gängigsten Python Webframeworks - Django, FastAPI und Flask -
generell und am Beispiel einer zu entwickelnden Webanwendung bei dem Praxispartner "SARAD GmbH".\\
Das Ziel dieser Arbeit ist es, das optimale Framework für das Anwendungsbeispiel zu bestimmen.
Hierfür werden für die Softwarelösung wichtige Aspekte aufgestellt, wie zum Beispiel die
Komplexität, Performanz und Datenbank-Interaktion, und nach diesen recherchiert und gestestet. Nach
Recherche und Tests werden die Ergebnisse in einem gewichteten Vergleich gegenübergestellt und ein
optimales Framework bestimmt.\\Das Ergebnis dieser Arbeit zeigt, dass Flask, aufgrund der
minimalistischen, einfachen und modularen Grundherangehensweise, das optimale Framework für das
Anwendungsbeispiel ist und für die Erstellung der Software genutzt werden sollte.\\ Das am Ende
stehende Fazit erläutert zudem kurz die Eigenschaften der genannten Frameworks und verweist darauf,
dass das Ergebnis einer solchen Analyse stark in Abhängigkeit zum geplanten Softwareprodukt steht.

\clearpage
\tableofcontents
\clearpage
\listoffigures
\listoftables
\end{titlepage}

\section{Einleitung}\label{s:einleitung}
\paragraph{Motivation}
Mein Praxispartner, die SARAD GmbH, produziert als primären Unterneh-\\mensgegenstand Detektoren und
Melder für das radioaktive Isotop Radon, welche in verschiedenen Bereichen, wie z.B. dem Bergbau
oder Wohnhäusern, in Gebieten mit hoher Radon-Belastung, eingesetzt werden. Seit ca. 2 Jahren bietet
das Unternehmen die Dienstleistung an, dass die von den einzelnen Geräten gesandten Daten auf einem
zentralen MQTT-Broker gesendet und ebenso zentral verarbeitet werden.\\Diese Dienstleistung verlangt
eine sowohl ordentliche als auch übersichtliche Führung vorhandener Kundendaten, um das Erstellen
von Rechnungen zu vereinfachen. Derzeit werden diese Daten in einer unstrukturierten Datenbank, ohne
jegliche spezifische Visualisierung, gespeichert und die Rechnungszeiträume müssen händisch
angedacht werden. Als Teil meiner Arbeit bei der SARAD GmbH ist es meine Aufgabe eine zentrale
Webanwendung zu entwickeln, damit die zuständigen Mitarbeiter einen geordneten Punkt haben um alle
Daten visualisiert zu bekommen und eventuelle Verwaltungsaktionen durchzuführen.\\Als Student der
Medieninformatik ist dieses Thema sehr passend, da die Webentwicklung einer der großen, zentralen
Punkte der Medieninformatik darstellt. In diesem Projekt wird es mir also ermöglicht mein
Webentwicklungswissen zu vertiefen, Erfahrung für die Praxis zu sammeln und eventuell mein
zukünftiges Studium durch Praxiserfahrung zu unterstützen. Ebenso stellt das Projekt eine gute
Grundlage für die Erstellung einer Bachelorarbeit dar.
\paragraph{Problemstellung und Zielsetzung}
Nachdem eine Bestandsaufnahme getätigt wurde, ist der erste Schritt bei der Erstellung dieser
Webanwendung die Frage, welche Technologien ich verwenden möchte oder sollte. Eine zentrale Gruppe
der Technologien zur Erstellung von Webanwendungen sind die Webframeworks, wovon es eine sehr große
Anzahl in unterschiedlichen Programmiersprachen gibt. Aufgrund von Aspekten, auf die ich im Verlauf
dieser Arbeit eingehen werde, beschränkt sich die Betrachtung der Webframeworks auf jene, die für
die Programmiersprache Python entwickelt wurden. Trotz dieser Einschränkung, gibt es weiterhin eine
große Anzahl an verschiedenen Webframeworks. Da diese Anzahl den Rahmen dieser Arbeit sprengen
würde, werde ich mich auf die drei verbreitetsten Python-Webframeworks beschränken, wodurch sich die
Leitfrage dieser Arbeit ergibt:\begin{center}\textbf{Welches der Python-Webframeworks - Django,
FastAPI oderFlask - ist das geeignetste für dieses Projekt?}\end{center}Um diese Frage zu
beantworten, werde ich zu Beginn die zu Grunde liegenden Anforderungen auflisten und bereits
gefallene Entscheidungen erläutern, anschließend eine Übersicht zu jedem der drei Frameworks bieten
und diese grundlegend und an meinem konkreten Anwendungsbeispiel analysieren und vergleichen. Am
Ende steht eine auf Recherchen und den daraus resultierenden Analysen fundierte Entscheidung,
welches dieser drei Frameworks für die Entwicklung der oben genannten Webanwendung am besten
geeignet ist.

\section{Anforderungen}\label{s:anforderungen}
Im folgenden Abschnitt gehe ich auf die verschiedenen Anforderungsaspekte ein und lege den Grundsatz
für die Recherche, die Analyse und den Vergleich der Webframeworks.\\Wie bereits in der Einleitung
beschrieben stellt sich zuerst folgende Frage - Warum soll die Webanwendung mit der
Programmiersprache Python, anstatt mit einer etablierteren oder spezialisierteren Sprache, wie zum
Beispiel PHP oder JavaScript mit Node.js, entwickelt werden? Die Wahl der Programmiersprache
resultierte aus drei simplen, aber kritischen Gründen:
\paragraph{Eigene Vorkenntnisse}Ich habe bereits, durch vorherige berufliche und private Projekte,
insgesamt über 5 Jahre unterschiedlichste Erfahrungen in der Entwicklung mit Python gesammelt und
habe damit die meiste Erfahrung. Die Nutzung von Python reduziert somit die Entwicklungszeit immens,
da ich mich nicht in eine mir eher unbekannte oder gänzlich unbekannte Sprache einarbeiten muss.
\paragraph{Vorkenntnisse der Mitarbeiter}Die anderen Mitarbeiter der IT- und Entwicklungsabteilung
der SARAD GmbH haben ebenso Vorkenntnisse in der Nutzung von Python. Somit ist es für diese leichter
nachzuvollziehen, wie der Quellcode der Anwendung geschrieben ist und es ist ebenso einfacher
möglich bei Bedarf Änderungen zu tätigen, nachdem ich das Unternehmen verlassen habe.
\paragraph{Vorhandene Skripte}Neben der Datenvisualisierung sollen über die Anwendung ebenso Skripte
\footnote{die genannten Skripte werden für die Zertifizierung (Erstellung von Zertifikaten, Rückzug
der Zertifikate und Bearbeitung von Berechtigungen) und die Erstellung der
Windows-Installationsdateien genutzt. Eine interne Bezeichnung des überstehenden Meta-Skripts ist
"MUST - MQTT User Service Tools"} auf einem internen Server ausgeführt werden. Aktuell sind diese
Skripte nur teilweise in Python geschrieben, was sich jedoch zukünftig ändern wird.

Die Wahl der Programmiersprache grenzt zwar bereits die Anzahl der Webframeworks immens ein, dennoch
bleibt eine Anzahl von mehr als 50 Webframeworks\footnote{als Grundlage dieser Schätzung wurde die
Anzahl an Python-Frameworks genutzt, die bei den TechEmpower-Benchmarks vorhanden sind.} übrig. Wie
bereits in der Einleitung beschrieben, habe ich mich auf die drei Webframeworks Django, FastAPI und
Flask begrenzt. Die Gründe hierfür sind:
\begin{itemize}
	\item\textbf{Verbreitung:} Django, FastAPI und Flask sind die aktuelle verbreitetsten
	Webframeworks für Python. Das macht die Arbeit mit diesen Frameworks sehr viel einfacher, da es
	eine große Anzahl an verschiedenen Quellen, Tutorials oder bereits beantwortete Fragen im
	Internet gibt. Ebenso gibt es eine Großzahl von unterschiedlichen Meinungen zu diesen
	Frameworks, was die Recherche und Evaluation erleichtert.
	\item\textbf{Karrierechancen:} Auch wenn möglicherweise kleinere, besser für meinen
	Anwendungszweck geeignete Webframework- Projekte im Python-Kosmos existieren, kann ich bei der
	Arbeit mit Django, FastAPI und Flask wichtige Arbeitserfahrung sammeln, welche mir zukünftig
	weitere Karrierechancen ermöglichen können.
	\item\textbf{Vorerfahrungen und Interesse:} In der Vergangenheit konnte ich bereits erste
	Erfahrungen mit dem Django-Framework sammeln. Im Falle von FastAPI und Flask habe ich zwar noch
	nicht mit diesen gearbeitet, jedoch ist mein Interesse an diesen beiden Frameworks hoch, da ich
	bereits Einblicke in Projekte, die mit diesen Frameworks entstanden sind, gewinnen konnte.
	\item\textbf{Varianz:} Die drei Kandidaten bieten eine hohe Varianz in ihrer Herangehensweise
	als Webframework und spiegeln so grundlegend unterschiedliche Archetypen mit unterschiedlichen
	Fokuspunkten und Philosophien wider.
\end{itemize}
Um konkrete Gesichtspunkte für eine Recherche und Analyse dieser Webframeworks zu erhalten, ist es
notwendig, klare Anforderungen an die Softwarelösung zu definieren.\\Aus den Gesprächen mit meinem
Betreuer, Herrn Michael Strey, und eigenen Überlegungen wurden folgende Ansprüche an die
Softwarelösung gesetzt - sortiert in absteigender Priorität:
\begin{itemize}
	\item\textbf{Visualisierung einer Datenbank:} Die notwendigen Kunden- und Vertragsdaten sollen
	in einer Datenbank gespeichert werden. Die Aufgabe der Software soll es sein diese, Datensätze
	vollständig und tabellarisch im Browser darzustellen. Aufgrund der Art der Daten ist eine
	relationale Datenbank notwendig\footnote{zum Zeitpunkt der Verfassung ist noch nicht festgelegt
	welche Datenbanklösung verwerdet wird. Tendenziell wird entweder MariaDB oder SQLite verwendet}.
	Das bedeutet, dass die Software eine Art des ORM\footnote{Object-relational mapping} benötigt.
	\item\textbf{Komplexität:} Da die Softwarelösung rein intern eingesetzt werden soll, soll sie
	nicht unnötig komplex und groß sein. Bestenfalls sollen nur die wirklich benötigten Funktionen
	eingebaut werden.
	\item\textbf{Ausführung von Skripten auf einer internen Maschine:} Es soll über die Webanwendung
	die Möglichkeit bestehen Python-Skripte auf einem internen Server auszuführen. Hierfür wird eine
	API-Schnittstelle benötigt.
	\item\textbf{Eigene Einschätzung:} Da ich der einzige Entwickler dieser Softwarelösung sein
	werde, ist meine eigene Einschätzung der Frameworks, wie passend sie mit ihren Eigenschaften für
	dieses Projekt sind, ein wichtiger Gesichtspunkt. Ebenso ist mein eigenes Interesse an der
	Arbeit mit den Frameworks relevant.
	\item\textbf{Performanz:} Die Menge der Nutzer dieser Software wird sich auf insgesamt maximal
	4-5 Personen belaufen. Dementsprechend ist davon auszugehen, dass nur etwa 1-2 Anfragen
	gleichzeitig stattfinden werden. Diese Tatsache macht den Aspekt der Performanz zwar deutlich
	unwichtiger als er sonst bei Webanwendungen wäre, jedoch ist es immer noch ein relevanter Punkt.
	\item\textbf{Erweiterbarkeit:} Da sich Software und die Anforderungen an diese andauernd
	verändern, soll es Möglichkeiten geben, die Webanwendung zukünftig einfach zu erweitern, falls
	sich neue Anforderungen an das Produkt ergeben sollten.
\end{itemize}

\section{Frameworks}
Im folgenden Abschnitt gehe ich auf die einzelnen Frameworks ein. Dabei werde ich auf allgemeine
Informationen zum Framework an sich, die grundlegende Philosophie sowie die Schritte zu einer
simplen "Hello World"-Website\footnote{einfachste anzunehmende Anwendung - die Software gibt einfach
nur den Text "Hello World" zurück} eingehen. Nach jedem Beispiel steht eine kleine Aufstellung von
Aspekten die für oder gegen das Framework sprechen \footnote{diese Aspekte sind aus eigener
Erfahrung entstanden. Es wurde jeweils etwa 5 Stunden mit jedem Framework gearbeitet um diese
Aspekte benennen zu können.}.
\subsection{Django}
Django, ursprünglich entwickelt von Adrian Holovaty und Simon Willison, wurde im Juli 2005
veröffentlicht und ist damit das älteste der drei Frameworks, die sich im Vergleich befinden.
Aktuell befindet sich Django in der Version 6.0\footnote{am 3. Dezember 2025 veröffentlicht} und
wird heutzutage von der Django Software Foundation\footnote{NonProfit-Organisation zur Entwicklung
und Instandhaltung von Django} auf dem neusten Stand gehalten. Das Framework ist vollständig
Open-Source und ist zu etwa 97\% in Python\footnote{siehe \url{https://github.com/django/django}}
geschrieben. Das Framework wird als \textit{"The Webframework for perfectionists
with deadlines"} beschrieben und verfolgt als Full-Stack-Lösung eine
\textit{"batteries included"}-Philosophie, was so viel bedeutet, dass Django grundsätzlich mit allen
Modulen und Funktionen, die eine Webanwendung benötigt, kommt. Unter diesen Modulen und Funktionen
befinden sich zum Beispiel ein Entwicklungs-Webserver für Tests, ein Templating-System, ein
dynamisches Admin-Interface und ein Authentifikations-System. Zudem ist Django auch noch mit
Drittanbieter-Paketen erweiterbar.
\clearpage
\subsubsection{Anwendungsbeispiel}
Django liegt eine grundlegende Unterscheidung zu Grunde - \textbf{Projekte und Apps}.\\Eine App ist
eine vollständige Web-Applikation, die etwas ausführt. Apps können zum Beispiel eine Umfrage, ein
Blog oder die angeforderte Softwarelösung sein.\\Ein Projekt ist eine Sammlung von Konfigurationen
und Apps für eine bestimmte Website.\\Bei diesem System gilt, dass ein Projekt mehrere Apps
beinhalten kann und eine App in mehreren Projekten vorhanden sein kann.\\Im Falle dieses
Anwendungsbeispiels ist ein Projekt und eine App notwendig\footnote{selbiges würde, aus aktueller
Sicht, auch für die geplante Softwarelösung gelten.}.\\Nachfolgend werden die Schritte
\footnote{alle in dieser Arbeit beschrieben Aktionen wurden in einer Linux-Distribution
druchgeführt, falls man diese Schritte mit einem anderen Betriebssystem durchführen möchte, muss man
die entsprechenden Befehlsäquivalente benutzen} beschrieben um ein einfaches Django-Projekt mit
einer "Hello World"-App aufzusetzen:
\begin{enumerate}
	\item\textbf{Installation mittels pip}\footnote{Paketmanager von Python}\\Folgender Befehl wird
	ausgeführt:
	\begin{verbatim}
		python -m pip install Django
	\end{verbatim}
	\item\textbf{Projekt-Verzeichnis erstellen}\\Folgender Befehl wird ausgeführt:
	\begin{verbatim}
		mkdir django
	\end{verbatim}
	\item\textbf{Projekt initialisieren}\\Folgender Befehl wird ausgeführt:
	\begin{verbatim}
		django-admin startproject project django
	\end{verbatim}
	Nach diesem Shell-Befehl entsteht im \verb|django|-Verzeichnis folgende Dateistruktur:
	\begin{verbatim}
	django
	  manage.py
	  project/
	    __init__.py
	    asgi.py
	    settings.py
	    urls.py
	    wsgi.py
	\end{verbatim}
	\item\textbf{App erstellen}\\Zuerst muss das Arbeitsverzeichnis gewechselt werden:
	\begin{verbatim}
		cd django
	\end{verbatim}
	Dann muss folgender Befehl ausgeführt werden:
	\begin{verbatim}
		python manage.py startapp myapp
	\end{verbatim}
	Nach diesem shell-Befehl entsteht im \verb|django|-directory folgende zusätzliche\\
	Dateistruktur:
	\begin{verbatim}
		django
		  ...
		  myapp/
		    __init__.py
		    admin.py
		    apps.py
		    models.py
		    tests.py
		    views.py
		    migrations/
		      __init__.py
	\end{verbatim}
	\item\textbf{View erstellen und per URL aufrufbar machen}\\In der \verb|myapp/views.py|-Datei wird folgender Code eingefügt:
	\begin{verbatim}
		from django.http import HttpResponse

		def index(request):
    		return HttpResponse("Hello World")
	\end{verbatim}
	\clearpage
	Um diesen Inhalt per URL aufrufbar machen zu können muss eine entsprechende URL-Konfiguration erstellt werden. Diese sind in
	der Django-Umgebung Python-Dateien mit dem Namen \verb|urls.py|.\\In diesem Fall muss eine \verb|urls.py|-Datei im
	\verb|myapp/|-Verzeichnis erstellt werden.\\Sobald die Datei erstellt wurde, wird folgender Code in sie eingefügt:
	\begin{verbatim}
		from django.urls import path
		from . import views

		urlpatterns = [
    		path("", views.index, name="index")
		]
	\end{verbatim}
	Als letzten Schritt muss nun noch die URL-Konfiguration des Projektes an sich korrekt gesetzt werden. Hierfür schreibt man
	folgenden Code in die Datei\\\verb|django/project/urls.py|:
	\begin{verbatim}
		from django.contrib import admin
		from django.urls import include, path

		urlpatterns = [
			path("myapp", include("myapp.urls")),
			path("admin", admin.site.urls)
		]
	\end{verbatim}
	\item\textbf{Testaufruf}\\Um den Entwicklungsserver von Django zu starten, und somit die "Hello World"-Seite im Browser
	aufrufbar machen zu können, gibt man folgenden shell-Befehl im \verb|django|-Verzeichnis ein:
	\begin{verbatim}
		python manage.py runserver
	\end{verbatim}
	\clearpage
	Der Server ist nun hochgefahren und beim Aufruf der Seite \url{http://localhost:8000/myapp/} öffnet sich die oben stehende Website.
	\begin{figure}[t]
	\includegraphics[scale=0.7,page=1]{\pdf}\hfill
	\caption{Firefox Screenshot - Django "Hello World"-Seite}
	\label{f:graph1}\end{figure}
\end{enumerate}
\subsubsection{Pros und Contras}
Die offensichtlichen Vorteile von Django sind klar – \textbf{es wird sofort fast\footnote{eine
API-Schnittstelle fehlt und muss mit einem Community-Paket importiert werden - auf diesen Aspekt
wird im Abschnitt "Vergleich und Analyse" eingegangen} alles mitgeliefert}, was bei einer
Webapplikation benötigt werden könnte. Zudem wird eine \textbf{klare Dateistruktur} erstellt, die
bei jedem einzelnen Django-Projekt gleich ist, was speziell bei zukünftiger, erneuter Arbeit mit
Django ein Vorteil ist. \textbf{Die erstemalige Einarbeitung in diese Struktur gestaltet sich jedoch
als schwierig}. Zusätzlich zu der Struktur kommt der für Python-Verhältnisse \textbf{komplexe
Quellcode} und der \textbf{große Konfigurationsaufwand}. Ein großer Teil\footnote{geschätzt 40\% der
Zeit} bei der Einarbeitungszeit mit Django wurde dafür verwendet, in der Dokumentation
nachzuschlagen.
\clearpage
\subsection{FastAPI}
FastAPI, enwtickelt und instandgehalten von Sebastián "tiangolo" Ramírez, wurde im Dezember 2018
veröffentlicht und ist damit das jüngste der Framework-Kandidaten. Aktuell befindet sich FastAPI in
der Version 0.124.4\footnote{am 12. Dezember 2025 veröffentlicht} und ist ein Open-Source
API-Framework, welches zu 100\% in Python\footnote{siehe \url{https://github.com/fastapi/fastapi}}
geschrieben ist. FastAPI ist eine ausschließliche Backend-Lösung für APIs und hat somit keine
Frontendfunktionen, wie einen Template-Renderer. Im Kern und der Namensgebung von FastAPI befindet
sich die Geschwindigkeit des Frameworks, welche durch die grundlegend asynchrone Funktionabilität
entsteht.

\subsubsection{Anwendungsbeispiel}
\begin{enumerate}
	\item\textbf{FastAPI und Uvicorn installieren}\\FastAPI kann mittels \verb|pip| installiert
	werden. Da FastAPI keinen integrierten Entwicklungssever hat, wird der ASGI-Server "uvicorn"
	benötigt um FastAPI zu starten. Es wird folgender Befehl ausgeführt:
	\begin{verbatim}
		python pip install fastapi uvicorn
	\end{verbatim}
	\item\textbf{Projekt-Verzeichnis erstellen und Arbeitsverzeichnis wechseln}\\Folgende Befehle wird ausgeführt:
	\begin{verbatim}
		mkdir fastapi
		cd fastapi
	\end{verbatim}
	\item\textbf{Hello World-Code schreiben}\\Folgender Code wird in eine neu erstellte `main.py`-Datei geschrieben:
	\begin{verbatim}
		from fastapi import FastAPI

		app = FastAPI()

		@app.get("/")
		async def index():
    		return {"message": "Hello World"}
	\end{verbatim}
	\item\textbf{uvicorn-Server starten}\\Folgender Befehl wird ausgeführt:
	\begin{verbatim}uvicorn main:app --reload\end{verbatim}
	\clearpage
	Der Server ist nun hochgefahren und beim Aufruf der Seite \url{http://localhost:8000} öffnet sich oben stehende Website
	\begin{figure}[t]
	\includegraphics[scale=0.7,page=2]{\pdf}\hfill
	\caption{Firefox Screnshot - FastAPI "Hello World"-Seite}
	\note{Die von FastAPI erstellte "Hello World"-Seite ist tatsächlich ein JSON-Objekt und keine HTML-Datei, wie es normalerweise ist.}
	\label{f:graph2}\end{figure}
	
\end{enumerate}

\subsubsection{Pros und Contras}
FastAPI hat im direkten Vergleich zu Django nahezu keinen Konfigurationsaufwand und eine sehr kleine
Dateistruktur. Der Quellcode ist aufgrund meiner Vorerfahrungen mit Asynchronität in
Python\footnote{Vorerfahrungen entstanden hauptsächlich bei der Arbeit mit discord.py - eine
Schnittstelle zur Interaktion mit Bots auf der Kommunikationsplattform Discord} sehr leicht
verständlich. Auch die komplett automatisch erzeugte OpenAI \footnote{ehemals Swagger} Dokumentation
ist sehr nützlich. Es gibt jedochjedoch gibt es ein konkretes, großes Problem: FastAPI hat keinen
Frontend-Kontext, hierfür benötigte es weitere Frameworks, wie React.

\subsection{Flask}
Das dritte und letzte Framework in der Auswahl, Flask, wurde am 1. April 2010 ursprünglich als
Aprilscherz von Armin Ronacher veröffentlicht. Aktuell befindet sich Flask in der Version
3.1.2\footnote{am 19.08.2025 veröffentlicht} und wird heutzutage von der
"Pallets"-Organisation\footnote{Organisation, welche beliebte und weiter verbreitete
Python-Frameworks weiterentwickelt und instand hält, z.B. Jinja, ein Template-Framework} erweitert
und instand gehalten.
\clearpage
Das Framework ist zu 100\% in Python geschrieben undvollständig Open Source. Flask ist, wie Django,
eine Fullstack-Lösung, besitzt jedoch eine grundlegend gegensätzliche Philosophie –Minimalismus.
Flask wird als Micro-Framework bezeichnet und wird mit sehr wenigen Komponenten bereit gestellt. Das
Grundframework kann jedoch mit vielen verschiedenen Community-Paketen erweitert werden.

\subsubsection{Anwendungsbeispiel}
\begin{enumerate}
	\item\textbf{Flask installieren}\\Flask kann mittels \verb|pip| installiert werden. Es wird folgender Befehl ausgeführt:
	\begin{verbatim}
		python pip install Flask
	\end{verbatim}
	\item\textbf{Projektverzeichnis erstelln und Arbeitsverzeichnis wechseln}\\Folgende Befehle wird ausgeführt:
	\begin{verbatim}
		mkdir flask
		cd flask
	\end{verbatim}
	\item\textbf{Hello World-Code schreiben}\\Folgender Code wird in eine neu erstellte \verb|app.py|-Datei geschrieben:
	\begin{verbatim}
		from flask import Flask

		app = Flask(__name__)

		@app.route("/")
		def hello():
			return "Hello World"
	\end{verbatim}
	\item\textbf{Testserver starten}\\Es wird folgender Befehl ausgeführt:
	\begin{verbatim}
		flask run
	\end{verbatim}
	\clearpage
	Der Server ist nun hochgefahren und beim Aufruf der Seite \url{http://localhost:5000} öffnet sich folgende Website:
	\begin{figure}[t]
	\includegraphics[scale=0.7,page=3]{\pdf}\hfill
	\caption{Firefox Screenshot - Flask "Hello World"-Seite}
	\label{f:graph3}\end{figure}
\end{enumerate}

\subsubsection{Pros und Contras}
Flask bietet die Simplizität von FastAPI mit dem Funktionsumfang von Django, wenn man sich die
entsprechenden Pakete hinzuzieht, die man für sein Projekt benötigt. Für die notwendigen
SQL-Operationen benötigt es zum Beispiel eine Erweiterung\footnote{genutzt wurde flask-sqlalchemy},
eine Vorgehensweise, die sich bei der Arbeit mit Flask als roter Faden durchzieht. Ein potentielles
Problem ist, dass viele Erweiterungen von Flask von unterschiedlichsten Entwicklern erstellt werden,
weshalb eventuell Konflikte entstehen können. Diese eventuelle Problematik zeigte sich in meiner
Einarbeitung jedoch nicht.

\section{Analyse und Vergleich}
Im nachfolgenden Teil gehe ich auf jedes Kriterium genau ein und gebe jedem Framework in dem
Kriterium eine Punktzahl zwischen 0 und 5 Punkten, wobei 0 bedeutet, dass das Framework das
Kriterium gar nicht erfüllt und 5 bedeutet, dass das Kriterium bestens erfüllt wird. Jedes Kriterium
erhält ebenso eine tabellarische Übersicht zur Punkteverteilung mit Anmerkungen.

Für die konkrete Evaluation und den Vergleich der drei Webframeworks benötigt man Kriterien. Da
bereits in den Anforderungen die wichtigsten Gesichtspunkte grundlegend genannt wurden, müssen diese
jetzt konkretisiert werden.

\begin{table}[h!]
	\begin{center}
		\begin{tabular}{|c||c|c|}
			\hline
			Kriterium & Erläuterung & Wichtung\\
			\hline\hline
			Frontend & Hat dieses FW eine Frontend-Komponente & 2x\\
			\hline
			Datenbank-Integration & Beschreibt, ob das FW eine eingebaute & 2x\\
			& Datenbank-Integration hat oder wie schwer es &\\
			& ist eine zu integrieren &\\
			\hline
			Komplexität & Beschreibt, wie hoch die Komplexität des FW & 1,5x\\
			& ist und wie steil die Lernkurve ist um mit &\\
			& diesem FW zu arbeiten &\\
			\hline
			API-Schnittstelle & Beschreibt, inwieweit API-Schnittstellen mit & 1,5x\\
			& dem FW erstellt werden können und die &\\
			& Komplexität des Aufsetzens dahinter &\\
			\hline
			Eigene Einschätzung & Beschreibt, meine eigene Einschätzung des FW & 1x\\
			& in unterschiedlichen Aspekten, z.B. Präferenz, &\\
			& Vorerfahrungen und Interesse &\\
			\hline
			Performance & Beschreibt, die Performance des FW im Bezug & 1x\\
			& auf Anfragen und Inhaltsrendering &\\
			\hline
			Erweiterbarkeit & Beschreibt, wie viele Erweiterungs- & 1x\\
			& möglichkeiten das FW besitzt und wie es &\\
			& skaliert werden kann &\\
			\hline
			Community & Beschreibt, wie gut die Community etabliert & 1x\\
			& ist und wie leicht man Antworten auf seine &\\
			& Fragen bekommt &\\
			\hline
		\end{tabular}
		\caption{Evaluationskriterien und deren Wichtung}
	\end{center}
	\label{table:1}
\end{table}

\subsection{Frontend}
Da die Softwarelösung eine visuelle Webanwendung sein soll, ist eine Frontend-Kompo-\\nente absolut
unverzichtbar. Sowohl Django als auch Flask kommen in ihrer Standardfassung mit einer vollständigen
Frontend-Engine mit eingebautem Template-Rendering. Das bedeutet, dass HTML-Dateien als Templates
mit entsprechenden Variablen-Feldern gepseichert werden. Beim Aufruf der Website werden die
Template-Variablen mit entsprechenden Inhalten ersetzt und somit eine Website mit den entsprechenden
Daten angezeigt. Somit sind alle Anforderungen erfüllt und es werden 5 von 5 Punkten vergeben.

FastAPI ist von Grund auf ein reines API-Framework, besitzt also KEINE Frontend-Engine. Um
entsprechende Frontend-Inhalte zu produzieren würde es ein weiteres Framework benötigen. Da in
dieser Analyse ausschließlich die drei genannten Frameworks betrachtet werden, müssen 0 von 5
Punkten vergeben werden.

\begin{table}[h!]
	\begin{center}
		\begin{tabular}{|c||c|c|}
			\hline
			Framework & Bewertung & Anmerkung\\
			\hline\hline
			Django & 5/5 & Vollständige Frontend-Engine mit Template Rendering\\
			FastAPI & 0/5 & KEINE Frontend-Engine\\
			Flask & 5/5 & Vollständige Frontend-Engine mit Template Rendering\\
			\hline
		\end{tabular}
		\caption{Bewertungskriterium: Frontend}
	\end{center}
	\label{table:2}
\end{table}

\subsection{Datenbank-Integration}
Um die richtigen Inhalte an das Template-Rendering zu übergeben benötigt es eine entsprechende
Datenbank-Integration. Weiter oben wurde bereits festgelegt, dass aufgrund der Art der Daten eine
relationale Datenbank notwendig ist, was automatisch bedeutet, dass ein entsprechendes ORM vonnöten
ist.

Django hat out-of-the-box ein integriertes ORM, somit können volle 5 von 5 Punkten vergeben werden.
Sowohl FastAPI als auch Flask kommen nicht direkt mit einem integrierten ORM, jedoch gibt es für
beide Frameworks eine konkrete Lösung: Bei FastAPI wird empfohlen SQLModel zu nutzen, ein ORM,
welches vom selben Entwickler erschaffen wurde. Im Falle von Flask existiert ein Modul mit dem Namen
"Flask-SQLAlchemy", welches ein konkretes Flask-Modul des weit verbreiteten Python-Paket
"SQLAlchemy"\footnote{dieses Paket beinhält sowohl ein ORM, als auch ein Toolkit zur Arbeit mit
relationalen Daten} ist. Aufgrund der bereits etablierten Modularität von Flask habe ich mich dafür
entschieden Flask in diesem Punkt etwas höher als FastAPI zu bewerten. Dementsprechend erhält
FastAPI 3.5 von 5 Punkten und Flask 4 von 5 Punkten.
\begin{table}[h!]
	\begin{center}
		\begin{tabular}{|c||c|c|}
			\hline
			Framework & Bewertung & Anmerkung\\
			\hline\hline
			Django & 5/5 & integriertes ORM\\
			FastAPI & 3,5/5 & nicht integriert, aber gute Zusammenarbeit mit\\
			& & SQLModel\\
			Flask & 4/5 & nicht integriert, aber konkretes Flask-Modul für\\
			& & SQLAlchemy\\
			\hline
		\end{tabular}
		\caption{Bewertungskriterium: Datenbank-Integration}
	\end{center}
	\label{table:3}
\end{table}
\clearpage

\subsection{Komplexität}
Da ich mit keinem der drei Frameworks viel Erfahrungen sammeln konnte, ist der Aspekt der
Komplexität und Lernkurve relevant, damit ich schnellstmöglich ordentlich mit diesen arbeiten und
die Webanwendung entwickeln kann.

Django hat aufgrund der komplexen Dateistruktur, dem Konfigurationsaufwand und hohen Anzahl von
integrierten Modulen eine sehr hohe Einstiegshürde\footnote{die offizielle Django-Dokumentation
besitzt, abseits von How-To und FAQ-Punkten, ~270 Punkte und Unterpunkte}. Deswegen erhält Django 2
von 5 Punkten. Sowohl FastAPI als auch Flask haben einfache Dateistrukturen und wenige Module
"out-of-the-box", was den Einstieg sehr vereinfacht. Häufig wird als Einstiegsschwierigkeit bei
diesen beiden Frameworks das Arbeiten mit Asnychronität genannt - diesen Punkt werde ich jedoch
außen vor lassen, da ich bereits viele Erfahrungen mit dem Entwickeln mit Asnychronität gesammelt
habe und dies somit keine Hürde mehr darstellt. FastAPI erhält deswegen 5 von 5 Punkten. Im Fall von
Flask ziehe ich jedoch einen Punkt ab, da die Modulauswahl und -nutzung am Anfang für Probleme
sorgen kann. Somit erhält Flask 4 von 5 Punkten.
\begin{table}[h!]
	\begin{center}
		\begin{tabular}{|c||c|c|}
			\hline
			Framework & Bewertung & Anmerkung\\
			\hline\hline
			Django & 2/5 & Komplexe Dateistruktur, viele integrierte Module\\
			FastAPI & 5/5 & Einfache Dateistruktur, ausschließlich API\\
			& & dementsprechend wenige Module\\
			Flask & 4/5 & Einfache Dateistruktur, wenige Grundmodule, eventuelle\\
			& & Schwierigkeiten beim Heraussuchen der benötigten\\
			& & Module\\
			\hline
		\end{tabular}
		\caption{Bewertungskriterium: Komplexität}
	\end{center}
	\label{table:4}
\end{table}

\subsection{API-Schnittstelle}
Um die Ausführung der Skripte zu gewährleisten, ist es notwendig, per API-Call einen Befehl an die
entsprechende Maschine zu senden. Dementsprechend ist eine API-Schnitt-\\stelle notwendig, um das
durchzuführen.

Django hat kein integriertes API-Framework. Hierfür benötigt man das "Django REST framework",
welches man separat installieren muss, ebenso ist diese Software nicht vom offiziellen
Django-Entwicklerteam. Django erhält in diesem Aspekt 3 von 5 Punkten - einen Punkt Abzug dafür,
dass das REST-Framework nicht integriert ist und einen weiteren, da es fragwürdig ist, dass eine so
zentrale Funktion eines Frameworks, bei einer "batteries included"-Philosophie, fehlt. FastAPI ist
ein reines API-Framework, demenstprechend erhält es 5 von 5 Punkten. Flask benötigt für API-Calls
ein Modul,  "FlaskRESTful", welches nicht standardmäßig integriert ist. Flask erhält deswegen 4 von
5 Punkten. 

\begin{table}[h!]
	\begin{center}
		\begin{tabular}{|c||c|c|}
			\hline
			Framework & Bewertung & Anmerkung\\
			\hline\hline
			Django & 2/5 & Nicht integriert - "Django REST framework" benötigt\\
			FastAPI & 5/5 & Integriert\\
			Flask & 4/5 & Nicht integriert - "FlaskRESTful" benötigt\\
			\hline
		\end{tabular}
		\caption{Bewertungskriterium: API-Schnittstelle}
	\end{center}
	\label{table:5}
\end{table}

\subsection{Eigene Einschätzung}
Da ich selbst mit einem der Frameworks arbeiten muss, ist es ebenso relevant wie meine Einstellung
zu jedem ist.

Ich habe bereits in zwei kleineren Projekte mit Django gearbeitet, dennoch habe ich nicht nicht
ansatzweise ausreichend Erfahrungen sammeln können, um einen Vorteil zu haben. Die hohe Komplexität,
das hohe Konfigurationsmanagement und die riesige Anzahl an Modulen sorgt dafür, dass ich Django in
diesem Projekt als "Overkill" erachte. Das bedeutet, dass das, was Django bietet, schlichtweg zu
viel ist und nicht benötigt wird. Aufgrund dieser Aspekte erhält Django insgesamt 2 von 5 Punkten.

Die Entwicklung von FastAPI war über die letzten Jahre für mich immer sehr interessant zu verfolgen
– nichtzuletzt, da auch Freunde und Bekannte intesiv mit diesem Framework arbeiten. Die Performanz
und Simplizität dieses Frameworks lässt mich vermuten, dass es zukünfitg eine größere Rolle spielen
wird. Das größte Problem, das FastAPI hat, ist, dass es ausschließlich eine Backend-Engine ist. Für
dieses Projekt ist das sehr ungünstig, da ich der einzige Entwickler bin und ich somit noch ein
zweites Framework, z.B. React, benötigen würde, um den Frontend-Teil zu realisieren. FastAPI erhält
deswegen 3 von 5 Punkten.

Anfangs war Flask simplerweise ein dritter Kandidat, jedoch hat sich diese Betrachtung über die
Arbeit hinweg stark geändert. Die grundlegende Einarbeitung und Flask an sich waren sehr interessant
für mich und die grundlegende Modularität macht es einfach, Projekte nach Bedarf zu erweiteren und
nur das zu nutzen, was auch wirklich benötigt wird. Häufig wird die Performanz von Flask als
Schwäche genannt, da diese sehr abhängig von den genutzten Modulen ist. Da dieses Projekt nur wenige
Module nutzen wird, sollte das jedoch kein Problem darstellen. Flask erhält deswegen die vollen 5
Punkte.

\begin{table}[h!]
	\begin{center}
		\begin{tabular}{|c||c|c|}
			\hline
			Framework & Bewertung & Anmerkung\\
			\hline\hline
			Django & 2/5 & Vorerfahrungen (aber zu wenige), sehr komplex und\\
			& & anstrengend, geringes Interesse, Overkill\\
			FastAPI & 5/5 & sehr hohes Interesse und vermutete Zukunftsrelevanz,\\
			& & leider suboptimal für dieses Projekt\\
			Flask & 4/5 & hohes Interesse, einfache Einarbeitung, ausreichend\\
			& & leistungsfähig für dieses Projekt\\
			\hline
		\end{tabular}
		\caption{Bewertungskriterium: Eigene Einschätzung}
	\end{center}
	\label{table:6}
\end{table}

\subsection{Performanz}
Auch wenn Performanz in diesem Projekt eher eine untergeordnete Rolle spielt, ist sie einer der
besten Aspekte eines Frameworks, um konkrete Vergleiche zu machen. Die Benchmarks beziehe ich von
TechEmpower's Round 21\footnote{vom 19.07.2022}, da dies das letzte Benchmark von TechEmpower ist,
in dem alle drei Kandidaten vorhanden sind. Von jedem Testtypen werden die "Best"-Tabellen
herangezogen, insofern sie vorhanden sind:
\begin{itemize}
	\item bei den Tests "Multiple queries" und "Data updates" wird die "20-queries"-Tabelle genutzt
	\item beim "Cached queries"-Test gibt es keine verfügbaren Daten
	\item für Flask wurde "flask-raw" genutzt, da "flask" nicht in allen Tests angetreten ist
	\item für Django wurde "django" genutzt
	\item für FastAPI wurde "fastapi" genutzt
\end{itemize}
\begin{table}[h!]
	\begin{center}
		\begin{tabular}{|c||c|c|c|c|}
			\hline
			Test & Django & FastAPI & Flask & Gewinner\\
			\hline\hline
			JSON serialization & 73.479 & 154.061 & 91.225 & FastAPI\\
			Single query & 19.402 & 62.020 & 44.033 & FastAPI\\
			Multiple queries & 1.396 & 12.406 & 13.574 & Flask\\
			Fortunes & 15.038 & 50.608 & 36.243 & FastAPI\\
			Data updates & 754 & 5.803 & 8.813 & Flask\\
			Plaintext & 79.188 & 155.667 & 101.348 & FastAPI\\
			\hline
		\end{tabular}
		\caption{TechEmpower Round21 Benchmark Ergebnisse}
	\end{center}
	\label{table:7}
\end{table}

Anhand dieser Tabelle sieht man, dass FastAPI in fast allen Fällen mit Abstand die beste Perfomanz
bietet und Django überall die schlechteste, teilweise nur ~1/8 der Performanz von FastAPI. Die
Ergebnisse von Flask sind generell etwas skeptisch zu betrachten, da die Performanz von Flask sehr
stark davon abhängig ist, welche und wie viele Module benutzt werden. Es ist aber davon auszugehen,
dass Flask in vielen Fällen eine bessere Performanz bietet als Django. Ein weiterer, wichtiger, 
erformanzbezogener Gesichtspunkt ist, dass Django standardmäßig snychron arbeitet, FastAPI und Flask
jedoch asynchron, wodurch mehrere Anfragen gleichzeitig besser verarbeitet werden können.
Bezugnehmend auf alle diese Aspekte bekommt Django 2 von 5 Punkten, FastAPI 5 von 5 Punkten und
Flask 3,5 von 5 Punkten.

\begin{table}[h!]
	\begin{center}
		\begin{tabular}{|c||c|c|}
			\hline
			Framework & Bewertung & Anmerkung\\
			\hline\hline
			Django & 2/5 & gute Performanz, snychron\\
			FastAPI & 5/5 & sehr gute Perfomanz, asynchron\\
			Flask & 3,5/5 & Performanz abhängig von genutzten Modulen,\\
			& & asynchron\\
			\hline
		\end{tabular}
		\caption{Bewertungskriterium: Performanz}
	\end{center}
	\label{table:8}
\end{table}

\subsection{Erweiterbarkeit}
Zwar ist die Erweiter- und Skalierbarkeit in diesem Projekt nicht von großer Bedeutung, jedoch für
die allgemeine Betrachtung der Frameworks. Es ist wichtig zu erwähnen, dass eine vertikale
Skalierung\footnote{vertikale Skalierung kann man vereinfacht als Erhöhung Rechenleistung der
bereits vorhandenen Maschinen beschreiben} immer eine Grenze hat, die sich am aktuellen
Technologiestand orientiert und eine horizontale Skalierung\footnote{horizontale Skalierung kann man
vereinfacht als Erhöhung der Anzahl benutzter Maschinen beschreiben} keine\footnote{es gibt
theoretische Grenzen, z.B. gesamt verfügbare Ressourcen auf dem Planeten Erde, jedoch sind diese
praktisch nicht relevant und für ein Unternehmen unerreichbar} Grenzen besitzt, wenn man finanzielle
Aspekte ignoriert.

Django, mit der "batteries included"-Philosophie, kommt als Gesamtpaket. Das bedeutet, dass eine
Erweiterung in vielen Fällen nicht notwendig ist, jedoch gibt es einige Möglichkeiten zur
Erweiterung\footnote{siehe Abschnitt "API-Schnittstelle"}. Aufgrund der primär synchronen
Arbeitsweise von Django ist eine vertikale Skalierung eher weniger von Relevanz, da die Arbeitslast
nicht asynchron vom gleichen Prozess bearbeitet werden kann, sondern per Loadbalancer an
verschiedene Prozesse, Prozessoren oder gar Maschinen verteilt werden muss. Django erhält damit
3 von 5 Punkten. FastAPI hat ebenso einige Erweiterungsmöglichkeiten \footnote{siehe Abschnitt
"Datenbank-Integration}, kommt aber schnell an seine Limitationen\footnote{siehe Abschnitt
"Frontend"}. Bezüglich Skalierung ist es bei FastAPI, aufgrund der asynchronen Arbeitsweise,
deutlich leichter vertikal zu skalieren. Damit erhält FastAPI ebenso 3 von 5 Punkten. Flask hat als
Kern eine modulare Philosophie und ist damit sehr leicht nach Bedarf erweiterbar. Jedoch ist zu
beachten, dass die Anzahl und Art von Modulen die Leistung des Frameworks beeinflussen kann. Wie
bereits bei FastAPI, sorgt die asynchrone Arbeitsweise von Flask dafür, dass vertikale Skalierung
sinnvoller zu realisieren ist. Flask erhält damit 5 von 5 Punkten.

\begin{table}[h!]
	\begin{center}
		\begin{tabular}{|c||c|c|}
			\hline
			Framework & Bewertung & Anmerkung\\
			\hline\hline
			Django & 3/5 & einige Erweiterungen möglich\\
			FastAPI & 3/5 & einige Erweiterungen, jedoch schnell Limitationen,\\
			& & gut vertikal skalierbar\\
			Flask & 5/5 & sehr viele Erweiterungen möglich, \\
			& & können aber Performanz beeinflussen, gut vertikal\\
			& & skalierbar\\
			\hline
		\end{tabular}
		\caption{Bewertungskriterium: Erweiterbarkeit}
	\end{center}
	\label{table:9}
\end{table}

\subsection{Community}
Alle drei Frameworks haben eine große Community, viele Ressourcen und zahlreiche Möglichkeiten
Antworten auf seine Fragen zu bekommen. Daher erhalten alle 5 von 5 Punkten.

\subsection{Zusammenfassung}
Es folgt eine Übersicht über alle oben genannten Aspekte und die vergebenen Punkte:

\begin{table}[h!]
	\begin{center}
		\begin{tabular}{|c||c|c|c|}
			\hline
			Bewertungskriterium & Django & FastAPI & Flask\\
			\hline\hline
			Frontend & 10 & 0 & 10\\
			Datenbank-Integration & 10 & 7 & 8\\
			Komplexität & 3 & 7,5 & 6\\
			API-Schnittstelle & 4,5 & 7,5 & 6\\
			Eigene Einschätzung & 2 & 3 & 5\\
			Performanz & 2 & 5 & 3,5\\
			Erweiterbarkeit & 3 & 3 & 5\\
			Community & 5 & 5 & 5\\
			\hline
			Summe & 39,5 & 38 & 48,5\\
			\hline
		\end{tabular}
		\caption{Zusammenfassung Bewertungskriterien}
	\end{center}
	\label{table:10}
\end{table}
FastAPI belegt in diesem Vergleich den letzten Platz, was eindeutig auf die Notwendig-\\keit der
Frontend-Engine zurückzuführen ist, die FastAPI nicht bietet. Blendete man diesen Aspekt aus,
belegte FastAPI einen sehr starken zweiten Platz. Django ist in den Aspekten Frontend und
Datenbank-Integration die beste Wahl, jedoch fällt die fehlende integrierte API-Schnittstelle, die
Komplexität dieses Frameworks  und schwache Performanz stark negativ auf. Das alles sorgte dafür,
dass Django am sclhectesten bewertet wurde. Flask zeigt sich durch die immense Flexibilität und
Simplizität, die es bietet, als Sieger dieses Vergleichs und damit als das passendste Framework für
dieses Projekt.

\section{Fazit}
Alle verglichenen Frameworks haben, wie in dieser Arbeit ausführlich beschrieben, ihre Vor- und
Nachteile und eignen sich für unterschiedliche Projekte. Dementsprechend ist das Ergebnis nicht
allgemeingültig zu betrachten. \textbf{Django} ist aufgrund seiner grundlegenden Struktur und
beigelieferten Module eine gute Lösung für stabile, große Projekte, bei denen bestenfalls
verschiedene, selbst entwickelte, Module vermehrt eingesetzt werden können, um das volle Potential
der Apps-Projects-Unterscheidung auszuschöpfen, man benötigt jedoch viel Erfahrung, um mit diesem
Framework schnell und produktiv arbeiten zu können, anstatt regelmäßig in der Dokumentation
nachzulesen. \textbf{FastAPI} lädt dazu ein, eine klare Unterteilung von Entwicklern in Frontend und
Backend zu tätigen und legt einen klaren und deutlichen Fokus auf Geschwindigkeit und Performanz
der APIs. Es ist am besten für Projekte geeignet, bei denen diese Aspekte eine wichtige Rolle
darstellen. \textbf{Flask} ist am ehesten für kleinere Projekte geeginet, die nicht alle Module
benötigen, welche Django mitliefert. Das Framework hat ebenso eine niedrige Einstiegshürde und
eignet sich dementsprechnd für alle Entwickler, unabhängig vom Erfahrungsstand.

Im Fall des spezifizierten Softwareproduktes eignet sich, wie man aus der Analyse herauslesen kann,
Flask am deutlichsten, da sich die Vorteile von Flask mit den Anforderungen decken und die Nachteile
des Frameworks nahezu irrelevant sind.
\clearpage
\section{Quellenangabe}
Offizielle Dokumentationen:
\begin{itemize}
	\item Django: \url{https://docs.djangoproject.com/en/5.0/}\\(Letzter Aufruf: 08.01.2026 9:58)
	\item FastAPI: \url{https://fastapi.tiangolo.com/de/}\\(Letzter Aufruf: 08.01.2026 10:00)
	\item Flask: \url{https://flask.palletsprojects.com/en/stable/}\\(Letzter Aufruf: 09.01.2026 13:42)
\end{itemize}
Repositories
\begin{itemize}
	\item Django: \url{https://github.com/django/django}\\(Letzter Aufruf: 18.12.2025 16:53)
	\item FastAPI \url{https://github.com/fastapi/fastapi}\\(Letzter Aufruf: 19.12.2025 13:27)
	\item Flask \url{https://github.com/pallets/flask}\\(Letzer Aufruf: 07.01.2026 10:38)
\end{itemize}
Benchmarks:
\begin{itemize}
	\item TechEmpower: \url{https://www.techempower.com/benchmarks}\\(Letzter Aufruf: 05.01.2026 12:36)
\end{itemize}
Sonstige Quellen:
\begin{itemize}
	\item GeeksForGeeks: \url{https://www.geeksforgeeks.org/python/comparison-of-fastapi-with-django-and-flask/}\\(Letzter Aufruf: 08.12.2025 10)
	\item Medium - Kapil Khatik: \url{https://medium.com/@kapildevkhatik2/django-vs-flask-choosing-the-right-web-framework-b8857ae38c67}\\(Letzter
	Aufruf 04.12.2025 11:37)
	\item Medium - Sai Nitesh Palamakula: \url{https://medium.com/@sainitesh/flask-vs-django-vs-fastapi-a-beginners-guide-to-choosing-the-right-
	python-web-framework-e5dd888a3495}\\(Letzter Aufruf: 04.12.2025 10:33)
	\item Medium - Tom Harrison: \url{https://tomharrisonjr.medium.com/flask-vs-django-why-i-chose-django-d248066ddd45}\\(Letzter Aufruf: 04.12.2025
	10:59)
	\item Medium - Viraj Sharma: \url{https://medium.com/wetheitguys/django-vs-flask-comparison-and-basic-project-setup-0b8e37589b19}\\(Letzer Aufruf:
	04.12.2025 10:44)
	\item Medium - Anusha: \url{https://medium.com/@athatikonda12/%EF%B8%8F-flask-vs-django-vs-fastapi-pros-cons-use-cases-9d60ca078962}\\(Letzer
	Aufruf: 04.12.2025 10:21)
	\item Medium - Anas Issath: \url{https://medium.com/@anas-issath/6a8d9741e7ee?sk=b064e7d3f4e01746acf6f07d2d9a8838}\\(Letzter Aufruf: 04.12.2025
	11:09)
	\item Towards Data Science - David Moore: \url{https://towardsdatascience.com/fastapi-versus-flask-3f90adc9572f/}\\(Letzer Aufruf: 04.12.2025
	13:16)
	\item Wikipedia: \url{https://en.wikipedia.org/wiki/Django_(web_framework)}\\(Letzter Aufruf: 05.01.2026 11:18)
	\item Wikipedia: \url{https://en.wikipedia.org/wiki/Django_Software_Foundation}\\(Letzter Aufruf: 04.12.2025 13:11)
	\item Wikipedia: \url{https://en.wikipedia.org/wiki/FastAPI}\\(Letzter Aufruf: 05.01.2026 11:17)
	\item Wikipedia: \url{https://en.wikipedia.org/wiki/Flask_(web_framework)}\\(Letzer Aufruf: 05.01.2026 11:18)
\end{itemize}
\section{Hilfsmittel}
\begin{itemize}
	\item LaTeX-Vorlage: Pascal Michaillat - Minimalist LaTeX Template for Academic Papers: \url{https://github.com/pmichaillat/latex-paper}
	\item Comprehensive TeX Archive Network - The CTAN \url{http://ctan.org}
	\item Integrated Development Environment: Visual Sutdio Code 1.106
	\item Betriebssysteme: Windows 10 22H2, Windows 11 24H2, Ubuntu 24.04
	\item Genutzte LLMs: Google Gemini\footnote{Versionen 2.5 Flash, 3.0 Fast, 3.0 Pro}
	\item \url{dict.leo.org}
	\item \url{openthesaurus.de}
	\item \url{korrekturen.de}
	\item \url{silbentrennung24.de}
	\item \url{scholar.google.com}
	\item \url{google.com}
	\item Korrekturlesende: Dennis, Nicole
\end{itemize}

\end{document}